\section{Proof of the Bound \texorpdfstring{$\frac{5}{4}\Delta^2$}{5 Delta squared / 4}}\label{sec:uniform}

Let $G$ be a finite loopless multigraph with maximum degree $\Delta$.
We prove that $\omega(L(G)^2)\leq\frac{5}{4}\Delta^2$. We first
prove the bound for simple graphs. To extend it to multigraphs, we
count how many edges of a given strong clique join each pair of
adjacent vertices.

\subsection{The Bound \texorpdfstring{$\frac{5}{4}\Delta^2$}{5 Delta squared / 4} in Simple Graphs}
\label{sec:simple-proof}\label{sec:finite-model}\label{sec:identities}\label{sec:finite-instance}
\label{subsubsec:conditional-extension-counts}

By \textbf{Lemma~\ref{lem:regularization}}, a nonempty strong clique
in a simple graph of maximum degree $\Delta$ remains a strong clique
in a $\Delta$-regular simple graph. Let $G$ be this regular graph and
mark the set $S$ of edges in the given clique. We retain the other
edges of $G$, which may join the endpoints of disjoint marked edges.
The symmetric function
$c_G\colon V(G)\times V(G)\to\{0,1,2\}$ is given by
\[
 c_G(u,v)=
 \begin{cases}
  2,&uv\in S,\\
  1,&uv\in E(G)\setminus S,\\
  0,&u=v\text{ or }uv\notin E(G).
 \end{cases}
\]

\begin{defn}\label{def:patterns}
For an integer $n\geq2$, put $P_n=\{0,\ldots,n-1\}$. A \emph{pattern}
$F$ of order $n$ assigns a color $c_F(i,j)\in\{0,1,2\}$ to each pair
$ij=\{i,j\}$ of distinct positions in $P_n$, with $c_F(0,1)=2$. We
call $(0,1)$ the ordered root. The underlying graph $H$ of $F$ has
vertex set $P_n$ and edge set consisting of the pairs of color $1$
or $2$. We call $F$ \emph{connected} if $H$ is connected and
\emph{admissible} if the pairs of color $2$ form a clique in $L(H)^2$.
Equivalently, for every two disjoint pairs of color $2$, some pair
with one position in each has color $1$ or $2$.
Let $\mathcal F_n$ be the set of connected admissible patterns of order $n$.
\end{defn}

\begin{defn}\label{def:color-preserving-maps}
Let $F$ be a pattern of order $n$, and let $G$ be a finite simple graph
with marked set $S\subseteq E(G)$. A \emph{color-preserving map} from $F$
to $G$ is a map $f\colon P_n\to V(G)$ such that
$c_F(i,j)=c_G(f(i),f(j))$ for all distinct $i,j\in P_n$.
For fixed $G$ and $S$, let $h_F$ be the number of color-preserving maps
from $F$ to $G$.
\end{defn}

Distinct positions may map to the same vertex, in which case their
pair has color $0$.

\begin{lemma}\label{lem:pullback}
Let $G$ be a finite simple graph, let $S\subseteq E(G)$ be a clique
in $L(G)^2$, and let $n\geq 2$. Every map $f\colon P_n\to V(G)$ with
$c_G(f(0),f(1))=2$ gives an admissible pattern $F$ by setting
$c_F(i,j)=c_G(f(i),f(j))$ for all distinct $i,j\in P_n$.
\end{lemma}

\begin{proof}
Consider disjoint pairs $ij$ and $k\ell$ of color $2$ in $F$. Their
images are edges of $S$. If the image edges are disjoint, some
endpoint of one is adjacent to an endpoint of the other since $S$
is a clique in $L(G)^2$. The corresponding pair of positions has
color $1$ or $2$.

If the image edges have a common endpoint, we may name the positions
so that $f(i)=f(k)$. The pair $jk$ then maps to the marked edge
$f(j)f(i)$ and has color $2$. The same reasoning applies when the
image edges are equal.
\end{proof}

Let $G$ be a finite simple $\Delta$-regular graph with $\Delta>0$,
and let $S\subseteq E(G)$ be a nonempty clique in $L(G)^2$ of size $m$.
For $F\in\mathcal F_n$, define
\begin{equation}\label{eq:q-definition}
 q_F=\frac{h_F}{2m\Delta^{n-2}}.
\end{equation}
The factor $2m$ counts the possible images of the ordered root.
For the two-position pattern $E_\star$ with its unique pair marked,
we have $h_{E_\star}=2m$ and hence $q_{E_\star}=1$.

Two patterns $F,F'\in\mathcal F_n$ are equivalent if there is a
permutation $\pi$ of $P_n$ with
$c_F(i,j)=c_{F'}(\pi(i),\pi(j))$ for every pair $ij$ of distinct
positions. The preimage of $\{0,1\}$ under $\pi$ may be any pair of
color $2$ in $F$, with either ordering of its two positions. These
are isomorphism classes of the colored patterns with the root allowed
to move. We call them orbits, as in the data files.

For $k\geq2$, $F\in\mathcal F_k$ and $u\in P_k$, let
$\mathcal E(F,u)$ be the set of admissible patterns on $P_{k+1}$ whose
restriction to $P_k$ is $F$ and whose pair $uk$ has color $1$ or $2$.
Each such pattern is connected because its restriction to $P_k$ is
connected and the new position $k$ is adjacent to $u$.

\begin{lemma}\label{lem:coordinates}\label{lem:degree}
Suppose that $G$ is $\Delta$-regular and that its $m>0$ marked edges
form a strong clique. For every $k\geq2$ and $F\in\mathcal F_k$,
$q_F\in[0,1]$ and $q_{F'}=q_F$ for every $F'$ in the orbit of $F$.
For each $u\in P_k$,
\begin{equation}\label{eq:degree-extension}
 \sum_{F'\in\mathcal E(F,u)}q_{F'}=q_F.
\end{equation}
\end{lemma}

\begin{proof}
Since $F$ is connected, its underlying graph has a spanning tree
containing the root edge. Choose the image of the ordered root in
$2m$ ways. Then assign images to the remaining positions in order of
increasing distance from $\{0,1\}$ in the tree, so that each new
position has a tree neighbor whose image is already chosen. The pair
joining them has color $1$ or $2$, so the new image has at most
$\Delta$ choices. Thus $0\leq h_F\leq 2m\Delta^{n-2}$, and
$q_F\in[0,1]$ by \eqref{eq:q-definition}.

If $\pi$ is a permutation of $P_n$ with
$c_F(i,j)=c_{F'}(\pi(i),\pi(j))$ for all distinct $i,j$, then
$f\mapsto f\circ\pi$ maps the color-preserving maps of $F'$ to those
of $F$. Precomposing with $\pi^{-1}$ gives the inverse, so
$h_F=h_{F'}$. The two counts have the same denominator in
\eqref{eq:q-definition}, which gives $q_F=q_{F'}$.

Fix a color-preserving map $f$ of $F$. The image of $k$ can be any
of the $\Delta$ neighbors of $f(u)$, including neighbors already
used by $f$. Each choice determines the colors of all pairs involving
$k$ and therefore a unique extension pattern in $\mathcal E(F,u)$.
Admissibility follows from \textbf{Lemma~\ref{lem:pullback}}.
Conversely, a map counted by $h_{F'}$ for $F'\in\mathcal E(F,u)$
restricts to a color-preserving map of $F$ and sends $k$ to a neighbor
of the image of $u$. Thus
$\sum_{F'\in\mathcal E(F,u)}h_{F'}=\Delta h_F$.
Dividing by $2m\Delta^{k-1}$ gives \eqref{eq:degree-extension}.
\end{proof}

Index the $N=23{,}789$ relabeling orbits in
$\bigcup_{n=3}^6\mathcal F_n$ as in~\cite[Section~A.1]{FGSTSupplement}.
For each $j\in[N]$, set $q_j=q_F$ for any $F$ in orbit $j$.
We treat $q_{E_\star}=1$ as a constant. When we collect the terms of a
sum over labeled patterns by orbit, each labeled pattern contributes
once to its orbit coordinate.

The degree equations in $Aq=b$ are the identities in
\eqref{eq:degree-extension} written in the orbit coordinates.
In the row for $(F,u)$, each $F'\in\mathcal E(F,u)$ contributes
one to the column indexed by its orbit. If $F=E_\star$, the right
side of \eqref{eq:degree-extension} is $1$. Otherwise we subtract
one from the column indexed by the orbit of $F$ and set the right
side to $0$.

The pairs $(F,u)$ used for the degree equations are specified
in~\cite[Section~A.2]{FGSTSupplement}, with $F$ of order $2$, $3$,
$4$ or $5$. At order $5$, we use one pattern from each orbit and
take every position of that pattern in turn as $u$. The resulting
$3{,}516$ equations form the system
\begin{equation}\label{eq:linear-system}
 Aq=b.
\end{equation}
For example, the first row, for $F=E_\star$ and $u=1$, is
$q_1+q_2+q_3+2q_4+q_5=1$. The two extensions with color pairs
$(c_{F'}(0,2),c_{F'}(1,2))=(1,2)$ and $(2,1)$ contribute to
the same orbit, giving the coefficient of $q_4$.

A \emph{type} of order $k$ is a pattern $T\in\mathcal F_k$.
A \emph{flag of type $T$} is a pattern in $\mathcal F_n$, for some
$n\geq k$, whose restriction to $P_k$ is $T$. The positions of the
flag remain labeled when counting extensions. Let
$\mathcal L=(F_1,\ldots,F_s)$ be a list of flags of type $T$ of
orders $n_1,\ldots,n_s$, respectively. For a color-preserving map
$f$ of $T$ into $G$ and any flag $F$ of type $T$, let $h_F(f)$ be
the number of color-preserving maps from $F$ to $G$ that extend
$f$. For each $i\in[s]$, put $v_i(f)=\frac{h_{F_i}(f)}{\Delta^{n_i-k}}$.
Let $v(f)$ be the column vector with entries $v_i(f)$. Define the
\emph{moment matrix} for $T$ and $\mathcal L$ by
\begin{equation}\label{eq:moment-definition}
 M(T,\mathcal L)=\frac{1}{2m\Delta^{k-2}}
       \sum_{f:T\to G}v(f)v(f)^{\mathsf T}.
\end{equation}
The sum runs over all color-preserving maps of $T$ into $G$. If a
flag $F_i$ equals $T$, then $v_i(f)=1$ for each such map.

\begin{lemma}\label{lem:moment}
For any finite list $\mathcal L$ of flags of type $T\in\mathcal F_k$,
the moment matrix $M(T,\mathcal L)$ is positive semidefinite.
If $n_i+n_j-k\leq6$ for all $i,j\in[s]$, each entry is a nonnegative
integer linear combination of the coordinates of $q$ and the
constant $q_{E_\star}=1$.
\end{lemma}

\begin{proof}
By \eqref{eq:moment-definition}, for every $x\in\mathbb R^s$,
$x^{\mathsf T}M(T,\mathcal L)x$ is the sum of the squares
$(x^{\mathsf T}v(f))^2$ divided by $2m\Delta^{k-2}$, so it is
nonnegative.

To compute entry $(i,j)$, put $n=n_i+n_j-k$. In $P_n$, identify each position $r$ of $F_i$ with $r\in P_{n_i}$.
Identify position $r$ of $F_j$ with $r$ if $r<k$ and with $n_i+r-k$
otherwise. The flags agree on their intersection
$P_k$. The unspecified pairs have one position in
$P_{n_i}\setminus P_k$ and the other in $P_n\setminus P_{n_i}$.
Let $\mathcal U_{ij}$ be the set of admissible completions of the
colors of these pairs. Each completed pattern is connected because
the flags are connected and share the type. We write the product of
the two colored flags as the formal sum
\begin{equation}\label{eq:colored-flag-product}
 F_i\star_T F_j=\sum_{H\in\mathcal U_{ij}}H.
\end{equation}
The positions are labeled, so each admissible assignment of
the unspecified colors contributes one summand.

Fix a color-preserving map $f$ of $T$. A pair of extensions of $f$
to $F_i$ and $F_j$ gives a unique map on $P_n$, since the extensions
agree on $P_k$. The images of distinct positions may coincide,
including positions from different flags. Evaluating $c_G$ on the
images of each unspecified pair determines its color. The resulting
pattern belongs to $\mathcal U_{ij}$ by
\textbf{Lemma~\ref{lem:pullback}}. Conversely, a map counted by
$h_H(f)$ for $H\in\mathcal U_{ij}$ restricts to an extension of $f$
to each of the two flags. Counting these pairs and dividing by
$\Delta^{(n_i-k)+(n_j-k)}=\Delta^{n-k}$ gives
\begin{equation}\label{eq:colored-flag-evaluation}
 v_i(f)v_j(f)=\frac{\displaystyle\sum_{H\in\mathcal U_{ij}}h_H(f)}
                         {\Delta^{n-k}}.
\end{equation}
Thus the product in \eqref{eq:colored-flag-product} expresses the
product of the two normalized extension counts exactly.

Substitute \eqref{eq:colored-flag-evaluation} into
\eqref{eq:moment-definition}. Summing $h_H(f)$ over all maps
$f$ of $T$ gives $h_H$, since each color-preserving map of $H$
has a unique restriction to $T$. The total power of $\Delta$ in
the denominator is $(k-2)+(n-k)=n-2$. Hence
\eqref{eq:q-definition} gives
\begin{equation}\label{eq:moment-expansion}
 \bigl(M(T,\mathcal L)\bigr)_{ij}
 =\frac{\displaystyle\sum_{H\in\mathcal U_{ij}}h_H}{2m\Delta^{n-2}}
 =\sum_{H\in\mathcal U_{ij}}q_H.
\end{equation}
For the flags in \textbf{Figure~\ref{fig:flag-union}}, let $H_a$ be
the completion that assigns color $a$ to the pair $23$, for
$a\in\{0,1,2\}$. All three completions are admissible. In particular,
when $23$ has color $2$, the pairs $02$ and $13$ of color $1$ join
the marked pairs $01$ and $23$. The product in
\eqref{eq:colored-flag-product} is therefore $H_0+H_1+H_2$.
Both flags occur in the list of nine flags for the marked
two-position type in~\cite[Table~A.2]{FGSTSupplement}. Their product
gives the entry $q_{H_0}+q_{H_1}+q_{H_2}$ of the corresponding
moment matrix.

If $n=2$, the union is $E_\star$ and contributes the constant $1$.
If $3\leq n\leq6$, each summand contributes to the coordinate of
its relabeling orbit.
\end{proof}

In \eqref{eq:moment-expansion}, the coefficient of an orbit coordinate
counts the labeled completions in that orbit. If two flags differ
only by a permutation of their positions outside the type, that
permutation gives a bijection between their extensions of each
type map. We therefore use one representative of each class of
flags under permutations fixing the type positions.

The system uses $57$ pairs $(T_\ell,\mathcal L_\ell)$, with types of
order $2$, $3$ or $4$. The flag orders in
\cite[Table~A.2]{FGSTSupplement} ensure $n_i+n_j-k\leq6$ for every
pair of flags in the same list. For the graph counts, we set
$M_\ell(q)=M(T_\ell,\mathcal L_\ell)$. To write these matrices as in
\eqref{eq:intro-affine-gram}, let $(D_{\ell j})_{ab}$ count the labeled
completions in $\mathcal U_{ab}$ that belong to orbit $j$, where
$\mathcal U_{ab}$ is taken for the list $\mathcal L_\ell$.
A union of order $2$ occurs only when both flags equal the type
$E_\star$. This gives a $1$ at the first diagonal entry of $K_\ell$
in each of the two lists with this type, and every other entry of
every $K_\ell$ is $0$. By \textbf{Lemma~\ref{lem:moment}},
\begin{equation}\label{eq:positive-moments}
 M_\ell(q)\succeq0\qquad(\ell\in[57]).
\end{equation}

To express $\frac{m}{\Delta^2}$ in terms of $q$, we choose the
images of $(0,1)$ and $(2,3)$ independently from the $2m$ ordered
marked edges. Let $\mathcal B\subseteq\mathcal F_4$ consist of the
patterns in which both $01$ and $23$ have color $2$. For $j\in[N]$,
let $o_j$ be half the number of patterns in $\mathcal B$ belonging
to orbit $j$. Then
\begin{equation}\label{eq:objective-vector}
 o^{\mathsf T}q=\frac{1}{2}\sum_{F\in\mathcal B}q_F.
\end{equation}

\begin{lemma}\label{lem:finite-encoding}\label{lem:objective}
Let $G$ be a finite simple graph of maximum degree $\Delta>0$, and
let $S\subseteq E(G)$ be a strong clique of size $m>0$. For the system
specified in~\cite[Section~A]{FGSTSupplement}, there exists
$q\in[0,1]^N$ such that $Aq=b$, $M_\ell(q)\succeq0$ for all
$\ell\in[57]$, and
\begin{equation}\label{eq:objective-value}
 o^{\mathsf T}q=\frac{m}{\Delta^2}.
\end{equation}
\end{lemma}

\begin{proof}
Apply \textbf{Lemma~\ref{lem:regularization}} to $G$ with $D=\Delta$
and mark the $m$ edges of $S$ in the regular supergraph. Form $q$
from the counts in \eqref{eq:q-definition}. By
\textbf{Lemma~\ref{lem:coordinates}}, these counts lie in $[0,1]$
and depend only on their orbits. Equations \eqref{eq:linear-system}
and \eqref{eq:positive-moments} give $Aq=b$ and $M_\ell(q)\succeq0$
for all $\ell\in[57]$. It remains to evaluate $o^{\mathsf T}q$.

There are $(2m)^2$ choices for the two ordered marked images.
Each choice specifies a map of the four positions and therefore
an admissible pattern by \textbf{Lemma~\ref{lem:pullback}}.
Admissibility guarantees a pair of color $1$ or $2$ joining the
two marked pairs of positions. Thus the pattern is connected,
so the choice contributes to exactly one pattern in $\mathcal B$.
Conversely, a map counted by $h_F$ for $F\in\mathcal B$ sends both
pairs to ordered marked edges. Hence $\sum_{F\in\mathcal B}h_F=(2m)^2$.
Using this count in \eqref{eq:q-definition} and
\eqref{eq:objective-vector} gives
\[
 o^{\mathsf T}q
 =\frac{(2m)^2}{2\cdot2m\Delta^2}
 =\frac{m}{\Delta^2}.
\]
\end{proof}

It remains to bound the objective $o^{\mathsf T}q$. We give integer
coefficients for the bound in \textbf{Lemma~\ref{prop:certificate-principle}}
and verify them against the equations and flag products above.

The file \path{reproducibility/data/compact_certificate.json} stores
the vector $L\in\mathbb Z^{3516}$ of degree multiplier numerators.
For each block $\ell\in[57]$, it also gives a list of pairs
$(w_{\ell a},z_{\ell a})$ with $w_{\ell a}\in\mathbb Z_{\geq0}$ and
$z_{\ell a}\in\mathbb Z^{s_\ell}$, where $s_\ell$ is the number of
flags in $\mathcal L_\ell$. These pairs represent the weights
$\frac{w_{\ell a}}{24}$ and vectors $\frac{z_{\ell a}}{10{,}000}$
in a decomposition of $Q_\ell$. We use the common
denominator $\Lambda=24\cdot10{,}000^2=2{,}400{,}000{,}000$ and set
\begin{equation}\label{eq:integral-Q}
 B_\ell=\sum_a w_{\ell a}z_{\ell a}z_{\ell a}^{\mathsf T},
 \qquad Q_\ell=\frac{B_\ell}{\Lambda},
 \qquad \lambda=\frac{L}{\Lambda}.
\end{equation}
There are $4{,}110$ pairs with positive weight across the $57$ blocks.
The representation in \eqref{eq:integral-Q} shows that each $Q_\ell$
is positive semidefinite.

\begin{lemma}\label{prop:checked-certificate}\label{prop:quadratic-bound}\label{lem:rational-witness}
The numerators $L$ and $B_\ell$ in \eqref{eq:integral-Q} satisfy
\begin{align}
 &\rho=A^{\mathsf T}L-\sum_{\ell=1}^{57}D_\ell^*(B_\ell)
                  -\Lambda o\geq0,\label{eq:integer-residual}\\
 &b^{\mathsf T}L+\sum_{\ell=1}^{57}\tr(B_\ell K_\ell)
     =3{,}000{,}000{,}000=\frac{5}{4}\Lambda.
       \label{eq:integer-constant}
\end{align}
Consequently, every finite simple graph $G$ satisfies
$\omega(L(G)^2)\leq\frac{5}{4}\Delta(G)^2$.
\end{lemma}

\begin{proof}
The exact checks in~\cite[Section~B]{FGSTSupplement} verify the
degree and moment coefficients against \eqref{eq:degree-extension}
and \eqref{eq:moment-expansion}. With these coefficients, the
residual in \eqref{eq:integer-residual} has $23{,}590$ positive entries
and $199$ zero entries, and the constant in \eqref{eq:integer-constant}
is $3{,}000{,}000{,}000$. Thus $r=\frac{\rho}{\Lambda}\geq0$ and
$C=\frac{5}{4}$ for the multipliers in \eqref{eq:integral-Q}.
If $G$ has an edge, let $m$ be the size of a maximum clique in $L(G)^2$.
For the vector $q$ given by \textbf{Lemma~\ref{lem:finite-encoding}},
\textbf{Lemma~\ref{prop:certificate-principle}} now gives
$o^{\mathsf T}q\leq\frac{5}{4}$.
Since $o^{\mathsf T}q=\frac{m}{\Delta(G)^2}$ by \eqref{eq:objective-value},
we obtain $m\leq\frac{5}{4}\Delta(G)^2$.
If $G$ has no edges, the clique number is zero.
\end{proof}

We next record two consequences of the simple-graph count that we will
use in the proof for multigraphs. The first finds a marked induced
five-cycle when the clique is large enough. The second bounds the
degrees available around that cycle.

Let $G$ be a simple graph with maximum degree at most a positive
integer $D$, and let $m>0$ of its edges form a strong clique.
Mark these $m$ edges. Write $t=\frac{D}{2}$ and
$\mu=\frac{5}{4} D^2-m$. By \textbf{Lemma~\ref{prop:quadratic-bound}},
$\mu\geq0$. When $D$ is odd, $f(D)=5t^2-t+\frac{1}{4}$. A counterexample
with $m>f(D)$ therefore satisfies
\begin{equation}\label{eq:odd-counterexample}
 \mu\leq t-\frac{5}{4}.
\end{equation}

By \textbf{Lemma~\ref{lem:regularization}}, there is a $D$-regular
simple graph $\widehat G$ containing $G$ as an induced subgraph.
Only the original clique edges are marked. In particular every vertex
incident with a marked edge belongs to $G$.
\textbf{Lemma~\ref{lem:finite-encoding}} supplies the normalized
counts for $\widehat G$.

A marked induced five-cycle is an induced cycle of length five in
$\widehat G$ whose five edges are marked. Let $q_C$ be the normalized
count for the pattern with marked cycle $0,1,2,3,4,0$ and root $01$.
Every color-preserving map of this pattern is injective. Adjacent
positions have different images, and given two nonadjacent positions,
there is a third position adjacent to exactly one of them.
In the balanced blow-up with $b$ vertices in each part and all
edges marked, there are $10b^5$ such maps, from the ten cyclic
orderings and the choice of one vertex per part. Hence
$q_C=\frac{10b^5}{2(5b^2)(2b)^3}=\frac{1}{8}$ in that example.

Fix an ordered image $c_0,\ldots,c_4$ of the cycle, taking indices
modulo five. Let $U_i$ consist of the vertices whose neighbors on
the cycle are exactly $c_{i-1},c_{i+1}$ and whose edges to both
cycle vertices are marked. The sets $U_i$ are pairwise disjoint,
contain $c_i$ in $U_i$, and lie in $V(G)$ since every vertex in
them is incident with a marked edge.
\textbf{Figure~\ref{fig:odd-cycle-parts}} illustrates $U_0$. Define
\begin{equation}\label{eq:odd-core-defs}
 a_i=|U_i|,\qquad B=\bigcup_{i=0}^4 U_i,\qquad
 s_i=D-a_{i-1}-a_{i+1},\qquad R=\sum_{i=0}^4 s_i.
\end{equation}
By the definitions of the five sets, the neighbors of $c_i$ in $B$
are exactly $U_{i-1}\cup U_{i+1}$. Since $\widehat G$ is $D$-regular,
$s_i=|N_{\widehat G}(c_i)\setminus B|$. Summing over the five cycle
vertices gives $R=\sum_i|N_{\widehat G}(c_i)\setminus B|$, an upper
bound on the number of edges of $G$ from the cycle to $V(G)\setminus B$.

\begin{figure}[htbp]
\centering
\begin{tikzpicture}[every node/.style={font=\small}]
  \node[graph vertex,label=below:$c_0$] (c0) at (0,1.4) {};
  \node[graph vertex,label=left:$c_1$] (c1) at (-1.55,0.45) {};
  \node[graph vertex,label=below:$c_2$] (c2) at (-0.95,-1) {};
  \node[graph vertex,label=below:$c_3$] (c3) at (0.95,-1) {};
  \node[graph vertex,label=right:$c_4$] (c4) at (1.55,0.45) {};
  \node[graph vertex,label=above:$x$] (x) at (0,2.35) {};
  \draw[marked edge] (c0)--(c1)--(c2)--(c3)--(c4)--(c0);
  \draw[marked edge] (c1)--(x)--(c4);
\end{tikzpicture}
\caption{A vertex $x$ in $U_0$ and its two marked edges to the cycle.}
\label{fig:odd-cycle-parts}
\end{figure}


To express $R$ in terms of the counts, let $L_R$ be the following
linear form in the six-position coordinates. The restriction to the
first five positions is the marked induced cycle. The coefficient
of a labeled extension is the number of neighbors of the sixth
position on the cycle if that position either has an unmarked edge
to the cycle or is adjacent to two consecutive cycle positions.
All other extensions have coefficient zero. A vertex outside $B$
with a cycle neighbor satisfies at least one of these two tests.
Indeed, if it fails both, it has at most two cycle neighbors and
all edges to them are marked. With only one such neighbor $c_i$,
its marked edge to $c_i$ has no endpoint in common with the
opposite cycle edge $c_{i+2}c_{i+3}$ and no edge joins their endpoints.
With two nonconsecutive marked neighbors, it belongs to one of
the sets $U_i$. Conversely, vertices of $B$ fail both tests.

Thus the sum over the sixth position is $R$ for each ordered
cycle. The denominators in \eqref{eq:q-definition} for five and
six positions differ by a factor of $D$. If $q_C>0$, this gives
$L_R=\frac{q_C}{D}\mathbb E[R]$, where the expectation is over a
uniformly chosen color-preserving map of the cycle pattern. If $q_C=0$,
there is no such map and $L_R=0$.

The five-cycle certificates give
\begin{equation}\label{eq:odd-moment-bounds}
 \frac{\mu}{D^2}\geq\gamma_C\Bigl(\frac{1}{8}-q_C\Bigr)-\epsilon_C,
 \qquad
 \frac{\mu}{D^2}\geq\gamma_R L_R-\epsilon_R.
\end{equation}
Both errors $\epsilon_C$ and $\epsilon_R$ are nonnegative. For the
first inequality we use either
$\gamma_C\geq\frac{99999}{100000}$ and $\epsilon_C\leq2\cdot10^{-9}$,
or $\gamma_C\geq\frac{6449}{10000}$ and $\epsilon_C\leq10^{-9}$.
For the second we use $\gamma_R\geq\frac{56425}{100000}$ and
$\epsilon_R\leq10^{-9}$. The files and checks are specified in
\cite[Sections~B and~C]{FGSTSupplement}.

Suppose that $\mu\leq\mu_*$, $\gamma_C\geq g_C>0$ and
$\gamma_R\geq g_R>0$. Put
$q_*=\frac{1}{8}-\frac{\mu_*+\epsilon_C D^2}{g_C D^2}$.
The first inequality in \eqref{eq:odd-moment-bounds} gives
$q_C\geq q_*$ when $q_C<\frac{1}{8}$, and the bound holds directly
when $q_C\geq\frac{1}{8}$ because $q_*\leq\frac{1}{8}$. If $q_*>0$,
using the second inequality and $L_R=\frac{q_C}{D}\mathbb E[R]$
gives a cycle with $R\leq\mathbb E[R]$ and hence
\begin{equation}\label{eq:odd-choose-anchor}
 R\leq\frac{\mu_*+\epsilon_R D^2}{g_R Dq_*}.
\end{equation}
Verified upper bounds on either error can be used in this estimate.

\begin{lemma}\label{lem:simple-cycle-counts}
For the sets defined in \eqref{eq:odd-core-defs}, every vertex of
$U_i$ is adjacent to every vertex of $U_{i+1}$. Moreover,
$|B|=\frac{5D-R}{2}$ and $a_i=\frac{D-R}{2}+s_{i-2}+s_{i+2}$ for every $i$.
\end{lemma}

\begin{proof}
For every $x\in U_i$ and $y\in U_{i+1}$ we have $xy\in E(G)$. The
marked edges $xc_{i-1}$ and $yc_{i+2}$ are disjoint. The
definitions of the parts exclude $xc_{i+2}$ and $yc_{i-1}$, and
$c_{i-1}c_{i+2}$ is excluded because the cycle is induced. Since
the marked edges form a clique in $L(G)^2$, we must have
$xy\in E(G)$. Thus a vertex of $U_i$ has at most $s_i$ neighbors
outside $U_{i-1}\cup U_{i+1}$.

The sum $A=\sum_i a_i s_i$ bounds the number of incidences at vertices
of $B$ after excluding the edges between $U_i$ and $U_{i+1}$ for all
$i$. Each remaining edge with both ends in $B$ contributes two
incidences. Put $e_0=\sum_i a_i a_{i+1}$, the number of edges between
parts $U_i$ and $U_{i+1}$.
Solving \eqref{eq:odd-core-defs} gives
\begin{equation}\label{eq:odd-core-sums}
 a_i=t-\frac{R}{2}+s_{i-2}+s_{i+2},\qquad
 |B|=\frac{5D-R}{2}.
\end{equation}
Summing the five equations for $s_i$ gives the formula for $|B|$.
For $a_i$, the two disjoint pairs of indices
$\{i-2,i+1\}$ and $\{i-1,i+2\}$ are the neighbor pairs of
$i+2$ and $i-2$, respectively. Subtracting their sizes from $|B|$
gives $a_i=|B|-2D+s_{i-2}+s_{i+2}$.
\end{proof}

By substituting \eqref{eq:odd-core-sums} into $\sum_i a_i s_i$
and using $2e_0=\sum_i a_i(D-s_i)$, we obtain
\begin{equation}\label{eq:odd-core-capacity}
 A=tR-\frac{R^2}{2}+2\sum_i s_i s_{i+2},
 \qquad
 e_0=5t^2-\frac{Rt+A}{2}.
\end{equation}
When $D$ is odd, $R=5D-2|B|$ is odd as well. To bound $\sum_i s_i s_{i+2}$,
consider the cycle on the indices with edges $i,i+2$. If two
nonadjacent vertices have positive weights, move their combined
weight onto the vertex with the higher sum of neighboring weights.
With the other weights fixed, the sum of products does not decrease.
Repeating leaves positive weights on a clique of size at most two.
The sum of products is then at most $\frac{R^2}{4}$. When $D$ is odd, integrality reduces this bound to $\frac{R^2-1}{4}$, so
\eqref{eq:odd-core-capacity} gives $A\leq tR-\frac{1}{2}$.

\begin{lemma}\label{lem:odd-unit-slack}
If an induced five-cycle with five marked edges has $R=1$, then $D$ is odd and $m\leq f(D)$.
\end{lemma}

\begin{proof}
Let $i$ be the unique index with $s_i=1$. Every vertex of the other
four parts has degree $D$ from the consecutive pairs alone. Thus a
marked edge outside the consecutive pairs has both ends in
$U_i\cup(V(G)\setminus B)$. The marked edge $c_{i+2}c_{i+3}$ has both
ends in parts whose consecutive neighbors are outside $U_i$. Both
vertices already have degree $D$ from those pairs, so neither is
an endpoint of, or adjacent to an endpoint of, the proposed marked
edge. This contradicts the clique property. Hence $m\leq e_0$.
For $R=1$, \eqref{eq:odd-core-capacity} gives
$A=t-\frac{1}{2}$ and $e_0=5t^2-t+\frac{1}{4}=f(D)$.
\end{proof}


\subsection{The Bound \texorpdfstring{$\frac{5}{4}\Delta^2$}{5 Delta squared / 4} in Multigraphs}
\label{sec:multigraph}\label{sec:weighted-counts}

The simple support of a multigraph records every adjacency that can
join two marked edges. To count the marked edges themselves, we must
also record how many copies have each pair of endpoints.

Let $D\geq1$ be an integer. Take a loopless multigraph $G$ with
$\Delta(G)\leq D$ and fix a
nonempty strong clique of size $m$ in $G$. Mark the edges of the clique
and restrict the host to the subgraph induced by their endpoints.
Every edge joining two disjoint marked edges has both endpoints in this
subgraph. Keep all marked edges. Between each pair of adjacent vertices joined
only by unmarked edges, keep one of those edges. The simple support $G_{\rm s}$ of the reduced
host therefore retains all adjacencies needed for the marked edges to
form a strong clique. Color an edge of $G_{\rm s}$ red if its endpoints are joined by a
marked edge of $G$, and blue otherwise. Let $s$ be the number of red
edges. For each red edge $uv$, choose one marked edge with endpoints
$u,v$ as its \emph{base copy}. Call the remaining marked edges with
these endpoints \emph{extra copies}. The red edges form a strong
clique in $G_{\rm s}$.

Put
\[
 b=m-s,\qquad M=\frac{m}{D^2},\qquad S=\frac{s}{D^2},
 \qquad e=\frac{b}{m}.
\]
Here $M$ and $S$ are normalized clique sizes and $e$ is the fraction
of marked edges that are extra copies. Thus $S=M(1-e)$. For each edge $uv$ of $G_{\rm s}$ let $b_{uv}$ be its
number of extra copies, with $b_{uv}=0$ for blue edges. Write
$d_{\rm s}(v)=d_{G_{\rm s}}(v)$ and
$j_v=\sum_{u\in N_{G_{\rm s}}(v)}b_{uv}$.
Then $\sum_v j_v=2b$ and $d_{\rm s}(v)+j_v\le D$.

We now write $G$ for this reduced multigraph. To obtain a
$D$-regular multigraph, apply \textbf{Lemma~\ref{lem:regularization}} and give each new edge
color blue. The reduced host is induced in the completion, and all
new edges have multiplicity one. Thus the red edges, their
multiplicities and the adjacencies among their endpoints are
unchanged. In the next counts we use the simple support of the
completed multigraph.

Use colors $0,1,2$ for a pair with equal or nonadjacent images,
a blue edge and a red edge, respectively. Let $P$ be a connected colored pattern on $n$ positions
with at least one red edge. As in \textbf{Definition~\ref{def:patterns}},
any two red edges in $P$ either meet or have endpoints joined by a red
or blue edge. In the weighted counts, $P$ has no designated root edge.
Let $T$ be a set of red edges of $P$ that form a forest on the positions.
Define
\[
 Q(P,T)=\frac{1}{2mD^{n-2}}
       \sum_{\phi}\ \prod_{uv\in T} b_{\phi(u)\phi(v)},
 \qquad
 X(P,T)=\frac{Q(P,T)}{|\operatorname{Aut}(P,T)|}.
\]
Here $\phi$ runs over color-preserving maps of the positions of $P$
to the simple support, and $\operatorname{Aut}(P,T)$ is the group of
permutations of the positions that preserve the pair colors and the
edge set $T$. As before, color $0$ permits equal images. Each edge
of $T$ contributes its extra multiplicity, even if two edges of $T$
have the same image, as in \textbf{Figure~\ref{fig:weighted-maps}}.
The forest condition is imposed on the positions. We use patterns
with $2\le n\le6$ and $|T|\le3$. Write $X$ for the vector of coordinates
$X(P,T)$, using one representative for each isomorphism class of pairs
$(P,T)$.

\begin{figure}[htbp]
\centering
\begin{tikzpicture}[marked edge/.append style={draw=red!70!black}]
  \node[graph vertex,label=below:$a$] (a) at (0,0) {};
  \node[graph vertex,label=below:$b$] (b) at (1.05,0) {};
  \node[graph vertex,label=below:$c$] (c) at (2.1,0) {};
  \draw[marked edge] (a)--(b)--(c);
  \node[above=6pt] at (b) {$T=\{ab,bc\}$};
  \draw[-{Latex[length=2.2mm]}] (2.65,0)--node[above]{$\phi$}(3.65,0);
  \node[graph vertex,label=below:$x$] (x) at (4.25,0) {};
  \node[graph vertex,label=below:$y$] (y) at (5.55,0) {};
  \draw[marked edge] (x)--(y);
  \draw[marked edge] (x) to[bend left=25] (y);
  \draw[marked edge] (x) to[bend right=25] (y);
  \node[above=10pt] at (4.9,0) {$b_{xy}=2$};
\end{tikzpicture}
\caption{A weighted map with $\phi(a)=\phi(c)=x$ and $\phi(b)=y$.
Its weight is $b_{xy}^2=4$.}
\label{fig:weighted-maps}
\end{figure}


For the pattern $P$ consisting of a single red edge, summing over the
$2s$ oriented red edges of the simple support gives $Q(P,\varnothing)=1-e$. Summing their extra
multiplicities gives $Q(P,E(P))=e$, where $E(P)$ consists of the
single red pair. Each degree identity extends a fixed pattern
at a specified position. For a red extension, the unweighted and
weighted counts represent one base edge and the extra copies,
respectively. An extension to a blue neighbor counts its unique copy.
Together the blue and red extensions account for exactly
$D$ edge copies at the image of that position. The equations use
those extensions whose terms all occur in the model.

Choose two ordered marked edges independently, allowing the same
copy in both choices. Classify each choice as a base or an extra
copy. There are $(2s)^2$ ordered pairs of base copies,
$2(2s)(2b)$ ordered pairs with one base and one extra copy, and
$(2b)^2$ ordered pairs of extra copies. Each choice of the two
ordered copies determines the images of four positions in the
support. Their colors give a connected pattern because the two
edges belong to the strong clique. The forest weights count the
choices of an extra copy for each specified pair. As in
\eqref{eq:objective-vector}, divide the counts by $2mD^2$ and take
one half. The three contributions to the objective are
\[
 O_0=\frac{s^2}{mD^2},\qquad
 O_1=\frac{2sb}{mD^2},\qquad
 O_2=\frac{b^2}{mD^2}.
\]
Their sum is $M$ since $m=s+b$. More generally, adding an independent
ordered marked copy to a fixed pattern gives $2M Q(P,T)$
after normalization. Indeed, choosing the ordered copy multiplies the raw count by
$2m$, and the two new positions contribute a factor $D^2$ in the denominator. Replacing that copy by a base
or extra copy gives $2S Q(P,T)$ or $2Me Q(P,T)$. These formulas
give linear bounds when $M$, $S$ and $e$ are restricted to intervals.
For example, the bounds $e_-\le e\le e_+$ imply
\[
 O_2\le(e_-+e_+)\left(O_2+\frac{O_1}{2}\right)-e_-e_+M
\]
by expanding $M(e-e_-)(e_+-e)\ge0$.

Fix a type on $k$ positions and flags extending the type. In addition
to the colors, the type and flags specify forest edges as in $(P,T)$.
For a type map $\phi$, let $w(\phi)$ be the product of the extra
multiplicities on its specified forest edges. For the $i$th flag, on
$n_i$ positions, sum over the extensions of $\phi$ to that flag. Each
extension contributes the product of the extra multiplicities on its
specified forest edges outside the type. Let $v_i(\phi)$ be this sum
divided by $D^{n_i-k}$ and let $v(\phi)$ be the column vector with
entries $v_i(\phi)$. Then
\[
 \frac{1}{2mD^{k-2}}\sum_\phi w(\phi)v(\phi)v(\phi)^{\mathsf T}\succeq0.
\]
The matrix is a sum of outer products with nonnegative coefficients
because every extra multiplicity is a nonnegative integer. Its $(i,j)$ entry is a sum over pairs of extensions that agree on the type,
with denominator $2mD^{n_i+n_j-k-2}$. The two extensions may use the
same vertex outside the type. Group the pairs by the colors on the
remaining position pairs. Each resulting pattern uses the forest
edges in the type once, through $w(\phi)$. This writes each matrix
entry as a linear combination of coordinates $X(P,T)$. We use flag
lists for which these products lie in the specified pattern collection.
The same counting identity applies to the matrices formed from local
degrees and from counts of extensions satisfying degree constraints.

To bound $M$ we also need to count the edges at the vertices of an
induced red five-cycle in the reduced multigraph. Use the parts $U_i$
and remaining degrees $s_i$ from \eqref{eq:odd-core-defs} for the red
edges of its simple support. Call the vertices $c_0,\ldots,c_4$ of
the cycle its anchors, with indices modulo five, and put
$B=\bigcup_i U_i$. We call $B$ the core of this cycle.
Let $L$ be the number of blue pairs between $U_i$ and $U_{i+1}$,
summed over $i$. By \textbf{Lemma~\ref{lem:simple-cycle-counts}},
all pairs between these two parts are adjacent. Reserving one copy
for each such pair leaves at most $s_i$ incidences at each vertex of
$U_i$. Set $a_*=\frac{D-R}{2}$. Equations \eqref{eq:odd-core-sums}
and \eqref{eq:odd-core-capacity} applied to the simple support give
\begin{equation}\label{eq:multi-core-identities}
 a_i=a_*+s_{i-2}+s_{i+2},\qquad
 A=\sum_i a_is_i=\frac{RD}{2}-\frac{R^2}{2}
             +2\sum_i s_i s_{i+2}\le\frac{RD}{2}.
\end{equation}
In particular $R=5D-2|B|$ has the same parity as $D$.

We now bound the contribution of edges incident with vertices outside $B$.
Set $r=\frac{R}{D}$, $j=\frac{\sum_{i=0}^4j_{c_i}}{D}$,
$g=\frac{m}{D^2}-\frac{5}{4}$, and $\lambda=\frac{L}{D^2}$.

\begin{lemma}\label{lem:multigraph-cycle-counts}
If $g>0$, then $r>\frac{1}{2}$. For $\frac{1}{2}\le r\le1$,
\begin{equation}\label{eq:multi-joint}
 2g+2\lambda\le
 \frac{1+r}{2+r}\,r\left(r-\frac{1}{2}\right)
 -\frac{1+r^2}{2+r}\,j.
\end{equation}
\end{lemma}
\begin{proof}
The signature of a vertex $x$ is the set $\{i:xc_i\in E(G_{\rm s})\}$,
of size $k_x$. Summing the capacity $s_i$ at the five anchors gives
$K=\sum_{x\notin B}k_x\le R$. For $x\notin B$, let $h_x$ be its marked
degree, let $Q_x=|N_{G_{\rm s}}(x)\cap B|$, and let $b_x$ count its
blue pairs to $B$. For any marked edge $xy$, the union of the
signatures of $x$ and $y$ covers every edge of the cycle. In fact,
if $xy$ is disjoint from a cycle edge, the strong clique condition
gives an adjacency to one of the ends of that cycle edge. If $xy$
shares an endpoint with a cycle edge, its other cycle endpoint is
a neighbor of the shared vertex. Since a vertex cover of $C_5$ has
at least three vertices, a vertex with $k_x=0$ can have a marked
edge only to a vertex $y$ with $k_y\ge3$. In particular such an $x$
has no marked edge into $B$, since $k_y=2$ for $y\in B$.
A vertex outside $B$ of signature size two has at least one blue anchor pair. For
consecutive positions, a marked anchor pair would fail the strong
clique condition with the cycle edge opposite those positions. For
nonconsecutive positions, two marked anchor pairs would put the vertex
in $B$.

Every blue pair between consecutive classes uses two unmarked incidences
in $B$. The $b_x$ blue pairs of $x$ into $B$ use one unmarked incidence
in $B$ each, disjoint from the first set of incidences. The two
counts are illustrated in \textbf{Figure~\ref{fig:multigraph-cut}}.
Summing the marked incidences inside and outside $B$ gives
\begin{equation}\label{eq:multi-core-cut}
 2m\le D|B|-2L+\sum_{x\notin B}(h_x-b_x).
\end{equation}

\begin{figure}[htbp]
\centering
\begin{tikzpicture}[
  marked edge/.append style={draw=red!70!black},
  blue pair/.style={draw=blue!70!black,densely dashed,line width=0.95pt},
  font=\small
]
  \filldraw[fill=black!4,draw=black!40,rounded corners=4pt]
    (-0.65,-0.55) rectangle (4.15,1.85);
  \node[anchor=north west] at (-0.65,1.85) {$B$};
  \node at (0.15,1.51) {$U_i$};
  \node at (3.3,1.51) {$U_{i+1}$};
  \node[graph vertex] (a) at (0.15,0.95) {};
  \node[graph vertex] (b) at (0.15,0) {};
  \node[graph vertex] (c) at (3.3,0.95) {};
  \node[graph vertex] (d) at (3.3,0) {};
  \node[graph vertex,label=above:$x$] (x) at (5.95,0) {};
  \draw[marked edge] (a) to[bend left=9] (c);
  \draw[marked edge] (a) to[bend right=9] (c);
  \draw[marked edge] (a)--(d);
  \draw[marked edge] (b)--(c);
  \draw[blue pair] (b)--node[below] {two incidences} (d);
  \draw[blue pair] (d)--node[above] {one incidence} (x);
\end{tikzpicture}
\caption{A blue edge between consecutive parts contributes twice to
the degree sum on $B$. A blue edge from $x\notin B$ to $B$ contributes
once. Both copies of a marked pair contribute to marked degree.}
\label{fig:multigraph-cut}
\end{figure}


We first bound the total outside contribution after charging $b_x$ against
the core capacity. Suppose $k_x=1$ and $x$ has a marked edge to a vertex of $B$.
We show that $b_x\ge a_*+1$. Let $j$ be the sole index in the signature of $x$, and write $xy$ for
its marked edge with $y\in U_i$. Since $\{j,i-1,i+1\}$ covers the
cycle, $i=j\pm2$. Of the two cycle neighbors of $c_j$, let $c_k$
be the one nonadjacent to $y$. For each $u\in U_j$ the edge $c_ku$
is marked. Comparing it with $xy$ gives $xu$ or $yu$ in the support.
The edge $xy$ uses one of the $s_i$ extra incidences at $y$, so
$y$ has at most $s_i-1$ neighbors in $U_j$. At least
$a_j-s_i+1\ge a_*+1$ vertices of $U_j$ are therefore adjacent to
$x$. Each of these edges is blue, since $\{j,j-1,j+1\}$ does not
cover the cycle. If $xy$ is marked with $k_x=1$ and $k_y=2$, the union of their
signatures is a vertex cover of the cycle and hence has size three.
Whenever the union on an outside marked edge has size three, its two
omitted positions are nonconsecutive. Every red edge from either
omitted anchor into $B$ forces its other end to be adjacent to at least
one end of the outside marked edge. The core neighborhoods of those
ends therefore cover three classes. The other two classes have total
size at most $D$, so the three classes contain at least $|B|-D=D+a_*$
vertices. Each end uses one degree unit on the outside marked edge,
and so each sees at least $a_*+1$ vertices of $B$. If the end of
signature size one has no marked edge into $B$, these pairs are blue.

We redistribute the quantities $h_x-b_x$ over vertices outside $B$.
A vertex of signature size one with a marked core neighbor or a
marked neighbor of signature size two has $b_x\ge a_*+1$. Since
$h_x+b_x\le D$, it retains a contribution at most $R-2$. Any other
vertex of signature size one sends $\frac{1}{2}$ per marked copy to
its neighbors, all of which have signature size at least three. Its
remaining contribution is at most $\frac{D}{2}$. A vertex $x$ of
signature size zero with a marked neighbor of signature size three
has $b_x=Q_x\ge a_*$ by the comparison with the two omitted anchors.
It distributes its positive contribution uniformly over its marked
edges. Since $h_x\le D-Q_x$, the amount per copy is at most
$\frac{2R}{D+R}$. The other vertices of signature size zero use the
same rule, with at most one per copy and with recipients of
signature size at least four. We discard any negative contribution
left at a vertex of signature size zero.

A recipient of signature size three with a marked neighbor of
signature size zero also sees at least $a_*$ vertices of $B$, so it
has at most $\frac{D+R}{2}$ marked copies outside $B$. Multiplying
this bound by the larger of the two sending rates shows that it
receives at most $\max\{\frac{D+R}{4},R\}$. If it has no marked
neighbor of signature size zero, it receives at most $\frac{D}{2}$.
A recipient with at least four positions receives at most $D$.
Adding each recipient's own contribution, at most $D$, shows that
the total contribution at every outside vertex is at most
$k_x\max\{\frac{D}{2},R\}$ for $R\le D$. By \eqref{eq:multi-core-cut},
\begin{equation}\label{eq:multi-rough-discharge}
2m\le D|B|-2L+\max\left\{\frac{D}{2},R\right\}K.
\end{equation}

If $R\le \frac{D}{2}$, the right-hand side of \eqref{eq:multi-rough-discharge}
is at most $D|B|+\frac{D}{2}R=\frac{5}{2}D^2$.
Thus $g>0$ implies $r>\frac{1}{2}$.

Assume now that $\frac{1}{2}\le r\le1$. The $K$ adjacencies from the
five cycle vertices to the exterior and the $\sum_i j_{c_i}$ extra marked
incidences at those vertices use disjoint units of the remaining
degree $R$. Thus $K\le R-\sum_i j_{c_i}$.
Divide all vertex counts and degrees by $D$, and all
edge counts by $D^2$. Write $\eta=\frac{1-r}{2}$, $v=\frac{1+r}{2}$
and $\rho=\frac{2r}{1+r}$. An outside marked edge whose signatures
have union of size three has two ends each with normalized core
degree at least $\eta$. Each end has outside marked degree at most
$v$.

Let $x$ count the vertices of signature size one with a marked
neighbor in the core or a marked neighbor of signature size two.
Their total contribution in \eqref{eq:multi-core-cut}, in the
normalized units, is denoted by $C_1$. Each of these vertices has
at least $\eta$ blue pairs into the core, and $C_1\le rx$. Let
$p_2,p_3$ count the marked edges from these vertices to vertices of
signature size two and of signature size at least three, respectively.
Let $n_2$ count their neighbors of signature size two via a marked
edge. Also let $z_3,z_0$ count the vertices of signature size three
and of signature size zero, respectively, that meet a marked edge
joining these two groups. Denote the count of such edges by $p_{03}$.
The three groups counted by $n_2,z_3,z_0$ are disjoint. Each of their
vertices uses at least $\eta$ of the core capacity by the comparison
with the two omitted anchors, and $p_2\le v n_2$ and
$p_{03}\le v\min\{z_0,z_3\}$.

By \eqref{eq:multi-core-identities}, the total normalized capacity
of the core after reserving its consecutive pairs is at most
$\frac{r}{2}$. Deduct the core capacity used by the groups counted
by $x,n_2,z_3,z_0$. To bound $C_1$, add the exterior marked
counts $p_2+p_3$ and deduct the blue core pairs at the vertices
counted by $x$. This gives
$C_1\le\frac{r}{2}-2\eta x-\eta(n_2+z_3+z_0)+p_2+p_3$.
With $\theta=\frac{2r-1}{2+r}$ and $c=r-\theta=\frac{1+r^2}{2+r}$,
take weights $1-\theta,\theta$ on the two bounds for $C_1$ to get
\begin{equation}\label{eq:multi-joint-first}
 C_1\le cx+\frac{\theta r}{2}
     +\theta(p_2+p_3)-\theta\eta(n_2+z_3+z_0).
\end{equation}

For the other outside vertices, use the redistribution leading to
\eqref{eq:multi-rough-discharge}. A vertex of signature size two
retains at most one. Assign $\theta p_2-\theta\eta n_2$ in
\eqref{eq:multi-joint-first} to the group of such vertices. Since
$p_2\le v n_2$, its additional contribution is at most $\theta r$
per vertex, or at most $c=\frac{1}{2}+\frac{\theta r}{2}$ per signature
position after adding the original contribution.

For a vertex of signature size at least four, let $d_1$ be its
marked degree to the group counted by $x$. None of those edges
transfers any contribution to it, so its own and received
contributions sum to at most $2-d_1$. The corresponding bound for a
vertex of signature size three with no marked neighbor of signature
size zero is $\frac{3}{2}-\frac{d_1}{2}$. Adding the corresponding
$\theta d_1$ from $\theta p_3$ gives at most $c$ per signature
position in both cases, since $c\ge\frac{1}{2}$ and $\theta\le\frac{1}{3}$.

For the group counted by $z_3$, write $d_{1,3}$ for its marked
degree to the group counted by $x$. The vertices counted by $z_3$
have outside marked degree at most $v$ and receive at most
$\frac{1}{2}$ per marked copy, except for edges to the vertices
counted by $z_0$, which can transfer up to $\rho$ each. The total
after adding the corresponding terms of \eqref{eq:multi-joint-first}
is at most
\[
 \left(1+\frac{v}{2}\right)z_3
 -\left(\frac{1}{2}-\theta\right)d_{1,3}
 +\left(\rho-\frac{1}{2}\right)p_{03}-\theta\eta(z_3+z_0).
\]
Use $z_0\ge\frac{p_{03}}{v}$, then $p_{03}\le v z_3$, to
bound this by $(1+\rho v-2\theta\eta)z_3=3c z_3$. The second
substitution is valid because
$\rho-\frac{1}{2}-\frac{\theta\eta}{v}\ge0$ for $\frac{1}{2}\le r\le1$.
This accounts for all terms of \eqref{eq:multi-joint-first}.
The vertices of signature size one outside the group counted by
$x$ retain at most $\frac{1}{2}$ each. Thus the right side of
\eqref{eq:multi-core-cut}, divided by $D^2$, is at most
$\frac{5-r}{2}-2\lambda+c(r-j)+\frac{\theta r}{2}$.
Subtracting $\frac{5}{2}$ proves \eqref{eq:multi-joint}.

\end{proof}


Apply the simple-graph five-cycle counts to the red edges of
$G_{\rm s}$. Completing the simple support to a $D$-regular simple
graph adds only blue edges. All vertices of a red five-cycle and of
its five parts are incident with red edges. They therefore belong
to the old graph, which is induced in the completion. Thus the
completion preserves both the red five-cycles and their five
parts. Write $q_C$ for the number of ordered induced red five-cycles
divided by $2sD^3$ and $L_R$ for the sum of $\frac RD$ over these
cycles divided by the same quantity. If each such cycle satisfies
$R\geq\alpha D$, then summing this inequality gives $L_R\geq\alpha q_C$.

Applying the simple-graph certificate inequality to four choices
of objective gives
\begin{equation}\label{eq:support-projection-four}
\begin{aligned}
 S-0.99999q_C
   &\le \frac{5}{4}-\frac{0.99999}{8}+2\cdot10^{-9},\\
 S+0.5682L_R
   &\le \frac{5}{4}+10^{-9},\\
 S+0.4897(L_R-0.32q_C)
   &\le \frac{5}{4}-\frac{0.4897\cdot0.32}{8}+2\cdot10^{-9},\\
 S+0.45119\left(L_R-\frac{q_C}{2}\right)
   &\le \frac{5}{4}-\frac{0.45119}{16}+2\cdot10^{-9}.
\end{aligned}
\end{equation}
The coefficients and the error allowances in
\eqref{eq:support-projection-four} are rational.
\textbf{Appendix~\ref{app:multi-projection}} specifies the
four certificate checks. To see why the displayed coefficients
suffice, write each certified inequality as $S+\gamma Y\leq
\frac54+\varepsilon$, moving the multiple of $\frac18$ to the left.
For $Y\geq0$, replacing $\gamma$ by a smaller nonnegative coefficient
preserves the inequality. For $Y<0$ the bound follows from
$S\leq\frac54$.

For a given $\alpha>0$, substitute $L_R\geq\alpha q_C$ in
\eqref{eq:support-projection-four}. Each of the four inequalities
now has the form $S\leq a_i-b_iq_C$. Include $S\leq\frac54$ as
a fifth inequality, with $a_5=\frac54$ and $b_5=0$. Thus
\[
 S\leq S_{\max}(\alpha):=
 \max_{x\geq0}\min_{1\leq i\leq5}(a_i-b_ix),
 \qquad e\geq1-\frac{S_{\max}(\alpha)}{M}.
\]
The second line of \eqref{eq:support-projection-four} has
$b_2=0.5682\alpha>0$. The minimum of the five lines is continuous
and piecewise linear and tends to $-\infty$ as $x$ tends to infinity.
Its maximum is therefore attained at zero or at the nonnegative
intersection of two lines with different slopes. This gives a finite
rational calculation for $S_{\max}(\alpha)$.

These inequalities can also be imposed directly on $X$. Multiplying
each simple-graph inequality by $1-e=\frac{s}{m}$ changes the cycle
normalization from $2sD^3$ to $2mD^3$. Its term $(1-e)S$ is $O_0$,
and the same factor changes $L_R$ to the corresponding weighted
coordinate sum. All terms of the resulting inequality are linear
in the weighted counts.

Let $k\geq3$ and let $K$ be the vertex set of a red $K_k$ in the weighted host.
For each vertex $v$ put $a_v=|N_{G_{\rm s}}(v)\cap K|$ and let
$h_v$ and $d_{\rm blue}(v)$ count the marked and the unmarked edges
incident with $v$, respectively. Let $m_K$ count the marked edges with
both endpoints in $K$, with multiplicity, and put
$J_K=\sum_{v\in K}j_v$. For a marked edge $uv$, the set
$(N_{G_{\rm s}}(u)\cup N_{G_{\rm s}}(v))\cap K$ meets every edge
of the red clique. Indeed, a disjoint red edge in $K$ must have an
endpoint adjacent to $u$ or $v$, and the same holds when the red
edge shares an endpoint with $uv$. Every vertex cover of $K_k$ has
size at least $k-1$, so $a_u+a_v\geq k-1$. If $u\in K$, then
$a_u=k-1$ and $a_v\geq1$. If both endpoints lie in $K$, then
$a_u+a_v=2(k-1)$. Summing over all marked edges with multiplicity gives
\[
 (k-1)m+\sum_{v\in K}h_v+(k-3)m_K
 \le\sum_v a_vh_v
 \le D(kD-J_K)-\sum_{v\in K}d_{\rm blue}(v).
\]
For the upper bound, $\sum_v a_v=\sum_{v\in K}(D-j_v)=kD-J_K$
in the regular completion. Moreover, $a_vd_{\rm blue}(v)$ is at
least the number of blue neighbors of $v$ in $K$. Summing this
last inequality counts each blue incidence at $K$ and gives the
upper bound on $\sum_v a_vh_v$, since $h_v=D-d_{\rm blue}(v)$.
Using $\sum_{v\in K}(h_v+d_{\rm blue}(v))=kD$ now gives
\[
 (k-1)m\le kD(D-1)-DJ_K-(k-3)m_K.
\]
For $k=5$, the right side is at most $5D(D-1)-20$. Thus $4m\leq5D(D-1)-20$ if a red $K_5$ exists.
This is impossible when $m\geq\frac54D^2$. For odd $D$ it is also
impossible when $m=f(D)+1$, since then $4m=5D^2-2D+5$.
The coordinates of every pattern containing a red $K_5$ are
therefore zero in either case.

Fix a red edge $xy$. Let $U=N_{G_{\rm s}}(x)\setminus (N_{G_{\rm s}}(y)\cup\{y\})$
and let $W$ be the vertices adjacent to neither $x$ nor $y$. In the
host induced on $U\cup W$, mark only the previously marked edges
between $U$ and $W$. Edges within a part can still join endpoints of
these marked edges. For each marked edge $uv$ between $U$ and $W$, and
each red neighbor $z$ of $y$, the strong clique condition for $uv$
and $yz$ forces $uz$ or $vz$ to be a host edge. All such vertices $z$ lie
outside $U\cup W$, and different $z$ give different edges incident
with $u$ or $v$ that are absent from the induced host. Hence the Ore-degree $\sigma$ of the bipartite
marked subgraph in this induced host is at most $2D$ minus the
number of red neighbors of $y$.

If there is a marked edge between $U$ and $W$, the corresponding
upper bounds are
\[
 \frac{2973}{10000}\sigma^2\quad\text{for all copies},\qquad
 \frac{273323}{10^6}\sigma^2\quad\text{for base copies},\qquad
 \frac14\sigma^2\quad\text{for extra copies}.
\]
We give the counts that yield these bounds. In the induced host
let $H$ be the bipartite graph of marked edges between $U$ and $W$,
with their multiplicities. Suppose $H$ is nonempty, put $h=\Delta(H)$, and take
a vertex $v$ of marked degree $h$. Let $A=N_H(v)$,
$B'=V\setminus N[v]$ (where $V$ and $N[v]$ refer to the induced
host), $p=|A|$ and $\ell=\sigma-d(v)$, with $d(v)$ the degree in
the host. Every $a\in A$ has host degree at most $\ell$ since $va$
is marked. Let $T$ be the number of marked edges meeting $A$,
$P$ the number of unmarked edges from $A$ to $B'$, and
$F=H[B']$. Write $f=|E(F)|$ and $u=\Delta(F)$, with $u=0$ if
$F$ is edgeless. Every marked edge omitted by the counts $T$ and
$f$ has an endpoint in $N(v)\setminus A$. This set has at most
$d(v)-h$ vertices. Consequently
\[
 |E(H)|\leq h(\sigma-h)-h\ell+T+f,\qquad T+P\leq p\ell .
\]
For each $a\in A$, its neighbors in $B'$ meet every edge of $F$.
Since neither endpoint of that edge is adjacent to $v$, one must
be adjacent to $a$ to join the edge to $va$. Thus $f\leq u\ell$.
The same observation shows that the Ore-degree of the marked edges
in the host induced on $B'$ is at most $\sigma-p$. Indeed, for
each of the $p$ vertices of $A$, at least one edge to an endpoint
of the marked edge is lost when the host is restricted to $B'$.

First count all marked copies. On each pair joined by a marked
edge, remove the unmarked copies. On every other adjacent pair
retain one unmarked edge. This preserves all marked edges and all
adjacencies, and $\sigma$ still bounds the degree sum on each marked
edge. Use the counts above for this reduced host. Let $X'$ comprise the vertices of
$B'$ in the same part of $H$ as $v$. For $x\in X'$, let $h_x$
count the marked edges from $x$ to $A$, with multiplicity, and put
$f_x=d_F(x)$.
Each marked edge of $F$ has one endpoint in $X'$. For every
$a\in A$ it has an endpoint adjacent to $a$. The adjacency to $a$
either has an unmarked copy or is marked and meets the endpoint
in $X'$. Counting these possibilities yields
\[
 pf\leq uP+\sum_{x\in X'} h_xf_x .
\]
If $T,f>0$, let $\alpha=\frac{T}{T+f}$ and $\gamma=\frac{f}{T+f}$.
The identity
$(h_x+f_x)(\gamma^2h_x+\alpha^2 f_x)-h_x f_x
=(\gamma h_x-\alpha f_x)^2$ and the bound $h_x+f_x\leq h$
give $\sum_{x\in X'} h_xf_x\leq h(\gamma^2T+\alpha^2 f)
=\frac{hTf}{T+f}$. Substituting the resulting lower bound for
$P$ in $T+P\leq p\ell$ gives
\[
 |E(H)|\leq h(\sigma-h)+\left(1-\frac hp\right)T
       +\left(1-\frac hu\right)f
       +\frac{h^2Tf}{pu(T+f)} .
\]
If $f=0$, the previous counts give $|E(H)|\leq h(\sigma-h)$.
The induced host on $B'$ has fewer marked edges, since it omits
all marked edges at $v$. Applying bounds for this smaller graph with
Ore-degree at most $\sigma-p$ gives upper bounds for $f$.
The proof notes in the repository~\cite{FGSTCode} give further necessary
inequalities by comparing the marked neighborhoods in $A$ of
two vertices of $X'$. For example, if $x,x'\in X'$ are nonadjacent,
every vertex of $N_H(x)\cap A$ that is not adjacent to $x'$ must
be adjacent to every vertex of $N_H(x')\cap A$ that is not adjacent
to $x$. Otherwise two marked edges would have no edge joining their
endpoints. The rational interval check in \textbf{Appendix~\ref{app:ore-replay}}
uses these neighborhood inequalities and the bounds for the smaller
graph. It excludes $|E(H)|>\frac{2973}{10000}\sigma^2$ over the full
range of $\frac{h}{\sigma},\frac ph,\frac uh,\frac{T}{\sigma^2}$
and $\frac{f}{\sigma^2}$.

For the base copies, pass to the simple support, mark one edge on
each marked pair, and take the host induced on their endpoints.
This preserves every adjacency that joins two disjoint marked edges.
Let $s'$ be the number of marked edges in this simple host and let
$\sigma'\leq\sigma$ be their Ore-degree. Every vertex has a marked
neighbor, so its degree is at most $\sigma'-1$. For a connected
pattern $P$ on $k$ positions rooted at an ordered marked edge,
normalize its map count by $2s'(\sigma')^{k-2}$. The resulting
coordinate $q_P$ lies in $[0,1]$ by the spanning-tree bound and
vanishes if the marked edges of $P$ contain an odd cycle.
The count of two ordered marked edges in \eqref{eq:objective-value}
now gives the objective $\frac{s'}{(\sigma')^2}$.
The extension sum at one position is at most $q_P$. If $ij$ is a
marked edge of $P$, the two extension sums at $i$ and $j$ together
are at most $q_P$, since the sum of the two image degrees is at
most $\sigma'$. The moment matrices remain positive semidefinite
by \textbf{Lemma~\ref{lem:moment}}. The rational certificate in
\textbf{Appendix~\ref{app:ore-replay}} combines the extension inequalities
with nonnegative multipliers and the moment matrices with
positive semidefinite matrices. Each negative residual coefficient
is charged its absolute value, using $0\leq q_P\leq1$. The resulting
upper bound on $\frac{s'}{(\sigma')^2}$ is
$\frac{27332270490572378277}{10^{20}}<\frac{273323}{10^6}$.
This proves the asserted bound for base copies.

Finally, mark only the extra copies between $U$ and $W$, reserving
one base copy on each originally marked pair as an unmarked edge.
Every adjacent pair of vertices now has an unmarked edge joining it.
Apply the same counts with $H$ the graph of these extra copies.
The result is immediate if $H$ is empty. Otherwise, for
every pair consisting of $a\in A$ and an edge of $F$, the
joining adjacency has an unmarked copy, and each such copy serves
at most $u$ edges of $F$. Therefore $pf\leq uP$.
If $f=0$, the first two inequalities give $|E(H)|\leq
h(\sigma-h)$ because $p\leq h$. If $f>0$, then $u>0$, and
\[
 |E(H)|\leq h(\sigma-h)-(h-p)\ell+
          \left(1-\frac pu\right)f
 \leq h(\sigma-h)+\left(1-\frac hu\right)f
 \leq h(\sigma-h).
\]
Here the second inequality uses $f\leq u\ell$ and $p\leq h$,
and the last uses $u\leq h$. Finally $h(\sigma-h)\leq
\frac{\sigma^2}{4}$.

Multiplying any of the three bounds by the nonnegative forest
weight of the root and summing over root maps gives a linear
inequality for $X$. Expanding a neighbor count uses one further
position, and expanding its square or counting another edge uses
two.

Write $AX=(1,0,\ldots,0)^{\mathsf T}$ for the degree and sampling
equations, $NX\le0$ for the linear inequalities, and $G_i(X)\succeq0$
for the matrix inequalities. To bound $M$, take a vector
$\lambda$ with first entry $\lambda_0$, a nonnegative vector $\rho$, and matrices $P_i$
given as nonnegative sums of rational rank-one outer products.
Let $G_i^*$ be the adjoint defined by
$(G_i^*(P_i))^{\mathsf T}X=\tr(P_iG_i(X))$. The constraints give $\rho^{\mathsf T}NX\leq0$ and
$\tr(P_iG_i(X))\geq0$. Thus, for every feasible $X$,
\[
 \left(A^{\mathsf T}\lambda+N^{\mathsf T}\rho
        -\sum_iG_i^*(P_i)\right)^{\mathsf T}X
 \le\lambda_0 .
\]
Let $r$ be the coefficient vector on the left minus the coefficient
vector of $M$. If $0\le X_j\le U_j$ for each coordinate $j$, then
$M\le\lambda_0+\sum_{j:r_j<0}(-r_j)U_j$. For a coordinate $X(P,T)$,
an upper bound is
$\frac{1-e}{|\operatorname{Aut}(P,T)|}$ for $T=\varnothing$ and
$\frac{e}{|\operatorname{Aut}(P,T)|}$ otherwise. To obtain it, start at
a red edge, in $T$ if $T$ is nonempty, and extend along a
spanning tree containing $T$. Each new position contributes at
most $D$ in support count or extra multiplicity. The certificates
also use bounds on sums of coordinates of equal order and
equal forest size, obtained by counting all labeled spanning
trees containing the specified red root. Positive residuals
can be dropped since the coordinates are nonnegative. With these bounds the total correction is affine in $e$, so
its maximum over the allowed interval is attained at an endpoint. For an infeasibility certificate, the objective
is zero and the compensated right side is negative.


\begin{lemma}\label{prop:multi-uniform}
Every strong clique of size $m$ in a finite loopless multigraph of maximum
degree at most $D$ satisfies $m\le\frac{5D^2}{4}$. If $m>0$, $D>0$, and
$M\ge\frac{4967}{4000}$,
then $e<\frac{1}{4}$.
\end{lemma}

\begin{proof}
If $m=0$ the asserted bound holds. Otherwise $D>0$, and the
red edges form a strong clique of size $s$ in the simple support,
which has maximum degree at most $D$. Thus
\textbf{Lemma~\ref{prop:quadratic-bound}} gives $S\leq\frac{5}{4}$.
The five-position certificates in~\cite[Section~G.4]{FGSTSupplement}
exclude $M\geq\frac{4967}{4000}$ and $e\geq\frac14$ simultaneously,
giving the second assertion. They also give
$e<\frac{421}{2000}$ for $M\geq\frac54$ and $M<1.28492$ for
$M>\frac54$.

Suppose that $M>\frac54$. For an induced red five-cycle put
$r=\frac RD$ and $g=M-\frac54$. The degree estimates in
\textbf{Lemma~\ref{lem:multigraph-cycle-counts}} give $r>\frac12$.
When $r\leq1$, \eqref{eq:multi-rough-discharge} implies
$2g\leq r(r-\frac12)$, and \eqref{eq:multi-joint} gives the
stronger inequality
\[
 2g\leq\frac{1+r}{2+r}r\left(r-\frac12\right).
\]
Both right-hand sides are increasing on $[\frac12,1]$.
For a proposed lower bound $M_-$ on $M$, substitute
$g\geq M_--\frac54$ and solve for a lower bound $\alpha$ on $r$.
For example, if $M_-=1.283$, then
$0.6084(0.6084-\frac12)<2(1.283-\frac54)$, so $r>0.6084$.
All values of $\alpha$ in \textbf{Table~\ref{tab:uniform-chain}}
are at most one and hence also apply when $r>1$.
It follows from \eqref{eq:support-projection-four} that
$e\geq1-\frac{S_{\max}(\alpha)}{M_-}$. This argument includes the
case of no induced red five-cycle, when $L_R=q_C=0$.

\begin{table}[htbp]
\centering
\caption{Intervals excluded by the six-position certificates. Each row
gives a certified upper bound for $M$ strictly below $M_-$.}
\label{tab:uniform-chain}
\begin{tabular}{@{}rrrr@{}}
\toprule
$M_-$ & $M_+$ & $\alpha$ & Lower bound on $e$\\
\midrule
1.283 & 1.28492 & 0.6084 & 0.051\\
1.279 & 1.283 & 0.6446 & 0.049\\
1.274 & 1.279 & 0.6242 & 0.0445\\
1.268 & 1.274 & 0.5978 & 0.0395\\
$\frac54$ & 1.274 & 0.5 & 0.0225\\
\bottomrule
\end{tabular}
\end{table}

For each row, assume $M_-\leq M\leq M_+$. On this interval the
weighted counts satisfy the lower bound on $e$ in the last column
and $e<\frac{421}{2000}$. The rational certificates for these
intervals, including the corrections for negative residuals, are
specified in~\cite[Section~G.4]{FGSTSupplement}. The first three
rows successively exclude $M\geq1.283$, $M\geq1.279$ and
$M\geq1.274$. The last row now applies throughout the remaining
interval and gives
\[
 M\leq\frac{124968062205915704703985647}
                 {100000000000000000000000000}<\frac54,
\]
contrary to $M>\frac54$.
\end{proof}

The quadratic bound gives $m\leq f(D)$ for even $D$. For odd $D$
we shall need one further estimate on the number of extra copies.

\begin{lemma}\label{lem:multi-extra}
Suppose $b>0$ and set $A=\frac{\sum_v j_v^2}{2bD}$. Then
\[
 M\le 2-\left(2+\frac{1}{D}\right)A-\frac{1}{D^2},
 \qquad
 Me\le 2A(1-A)-\frac{A}{D}.
\]
In particular, $e\le 1-\frac{M}{2}$ whenever $M\ge1$.
\end{lemma}

\begin{proof}
Count ordered pairs $(g_1,g_2)$ of marked edges, allowing $g_1=g_2$.
For each ordered pair of marked edges, the strong clique property gives adjacent vertices
$u,v$ with $u$ an endpoint of $g_1$ and $v$ an endpoint of $g_2$.
If the edges share an endpoint, there are at least two
such ordered pairs $(u,v)$. Among the $b^2$ ordered pairs of extra copies,
exactly $\sum_v j_v^2-\sum_{uv\in E(G_{\rm s})}b_{uv}^2$ share an
endpoint. The second sum removes the
double count of pairs sharing both endpoints. The result is at least
$\frac{1}{2}\sum_v j_v^2$, since $\sum_v j_v^2\ge2\sum_{uv}b_{uv}^2$.
Counting the ordered adjacent endpoints gives
\[
 b^2+\frac{1}{2}\sum_v j_v^2
 \le 2\sum_{uv\in E(G_{\rm s})}j_uj_v
 \le \sum_v j_v^2(D-j_v).
\]
The last inequality follows from $2j_uj_v\le j_u^2+j_v^2$ and
$d_{\rm s}(v)\le D-j_v$. The quantities $\frac{j_v}{D}$ have
weighted mean $A$ with weights $\frac{j_v}{2b}$. Expanding
$\sum_v \frac{j_v}{2b}\bigl(\frac{j_v}{D}-A\bigr)^2\geq0$
gives $\sum_v j_v^2(D-j_v)\le 2bD^2A(1-A)$.
This gives the upper bound on $Me=\frac{b}{D^2}$.

Let $h_v$ be the number of base copies incident to $v$. Among ordered pairs
with the extra copy first and the base copy second, at least
$b$ share an endpoint, since each extra copy shares its two
endpoints with its own base copy. Using $h_v\le D-j_v$ gives
\[
 sb+b\le \sum_v j_v\sum_{u\in N_{G_{\rm s}}(v)} h_u
 \le D\sum_v j_v d_{\rm s}(v)-2\sum_{uv\in E(G_{\rm s})}j_uj_v.
\]
Using $d_{\rm s}(v)\le D-j_v$ and the lower bound on
$2\sum_{uv\in E(G_{\rm s})}j_uj_v$ gives
\[
 b(m+1)\le 2bD^2-\left(D+\frac{1}{2}\right)\sum_v j_v^2.
\]
Divide by $bD^2$ to obtain the first inequality. For $M\ge1$ we
have $0\le A\le 1-\frac{M}{2}\le\frac{1}{2}$. The function $2x(1-x)$ is
increasing on $[0,\frac{1}{2}]$, so $Me\le M(1-\frac{M}{2})$.
\end{proof}

\begin{lemma}\label{lem:multi-even-equality}
Suppose the host has positive even maximum degree $D$ and the
number of marked copies is $\frac{5D^2}{4}$. The host induced by
the marked endpoints is a blow-up of $C_5$ with all five parts of
size $\frac{D}{2}$. Every edge of that blow-up is single and marked.
\end{lemma}

\begin{proof}
The five-position multiplicity bound applies at $M=\frac{5}{4}$ and
gives $e<0.2105$. If every induced red five-cycle had
$\frac{R}{D}\ge\frac{1}{2}$, the support bound in \eqref{eq:support-projection-four} would give $S\le1.22181$ and
hence $e>0.0225$. These bounds satisfy the hypotheses of the last six-position certificate in \textbf{Table~\ref{tab:uniform-chain}}, whose upper bound is strictly less than
$\frac{5}{4}$. Thus there is a cycle with $R<\frac{D}{2}$.

Apply the redistribution leading to \eqref{eq:multi-rough-discharge}
at this cycle. For every outside vertex $x$ with $k_x>0$, its
resulting contribution is strictly less than $\frac{k_xD}{2}$.
For $k_x=1$, a vertex either keeps at most $R-2$ or sends half of
its marked degree and has a blue anchor edge. For $k_x=2$, there
is a blue anchor edge and no incoming transfer. For $k_x=3$ with
a marked neighbor of signature size zero, the incoming transfer
is at most $\max\{\frac{D+R}{4},R\}<\frac{D}{2}$. Without such a
neighbor, it is at most $\frac{D-3}{2}$, since the three anchor
adjacencies use degree units inside the core. For $k_x\ge4$,
the transfer is at most $D-k_x$ and its own contribution is
at most $D$. Vertices of signature size zero retain no positive
contribution. Since $K\le R$, the presence of an outside vertex
of positive signature size would give $2m<D|B|+\frac{DK}{2}
\le\frac{5D^2}{2}$ by \eqref{eq:multi-core-cut}.

Hence $K=0$. Now $2m\le D|B|=\frac{D(5D-R)}{2}$, so $R=0$.
By \eqref{eq:multi-core-identities}, the five part sizes are
$\frac{D}{2}$. For a vertex of $U_i$, the neighbors in
$U_{i-1}\cup U_{i+1}$ already use all $D$ degree units. There are
$\frac{5D^2}{4}$ edges on these pairs. Any marked copy outside
$B$ would require an outside vertex of positive signature size.
Thus all edges in the blow-up are marked, and $B$ contains every
marked endpoint.
\end{proof}


We finish with a consequence used in \textbf{Section~\ref{sec:exact}}
to force an induced red five-cycle. For $m,D>0$ we have
\begin{equation}\label{eq:absence-red-cycle}
 q_C=0 \quad\Longrightarrow\quad M<\frac{98}{81}.
\end{equation}
Suppose that $q_C=0$ and $M\geq\frac{98}{81}$. The uniform bounds for
the multigraph and its simple support give $M,S\leq\frac54$.
If $b=0$ then $e=0$. If $b>0$, \textbf{Lemma~\ref{lem:multi-extra}}
gives $e\leq1-\frac{M}{2}$. Thus in either case
\[
 \frac{98}{81}\leq M\leq\frac{12883}{10000},
 \qquad 0\leq e\leq\frac{32}{81}.
\]
These are the intervals used in the stored calculation.

Let $P$ be the five-position pattern of an induced red five-cycle.
As in the definition of $q_C$, every color-preserving map of $P$ is injective. Hence $q_C=0$ excludes every such map. It follows from
the definition of $X(P,T)$ in
\textbf{Subsection~\ref{sec:weighted-counts}} that $X(P,T)=0$ for
each forest $T\subseteq E(P)$ used in the calculation, including
$T=\varnothing$. Since the coordinates are nonnegative, the verifier
imposes all these equations by summing the distinct coordinates
$X(P,T)$ and requiring this sum to be at most zero.

Together with this inequality, the five-position calculation uses
the weighted degree and flag constraints, the simple-support bound,
the first inequality in \eqref{eq:support-projection-four}, and the
rooted five-cycle bounds indexed in the proof notes accompanying~\cite{FGSTSupplement}.
Using the five-position certificate from Section~G.4 of
the supplement~\cite{FGSTSupplement} and adding the correction for
negative residual coefficients gives
\[
 M\leq\frac{14802480264533669793521}{12250000000000000000000}
   <\frac{98}{81}.
\]
This contradicts the assumed range of $M$ and proves
\eqref{eq:absence-red-cycle}.

